% IEEE Quantum Week (QCE) technical paper format: full paper 8-10 pages including figures,
% tables and appendices, plus up to 2 pages of references.
\documentclass[10pt,conference]{IEEEtran}
\usepackage{amsmath, amssymb}
\usepackage{braket}
\usepackage{booktabs}
\usepackage{graphicx}
\usepackage{xcolor}
\usepackage[hidelinks]{hyperref}

% An unfilled section prints FILL, which the /techne:latex draft-marker gate blocks.
\newcommand{\fillin}[1]{\textcolor{red}{\textbf{FILL:} #1}}

\begin{document}

\title{When Do Variational Quantum Circuits\\ Keep Federated Training Data Private?}

\author{\IEEEauthorblockN{AJ Barea}
\IEEEauthorblockA{Rochester Institute of Technology\\
Rochester, NY, USA\\
ajb6289@rit.edu}}

\maketitle

\begin{abstract}
\fillin{Question, method, the regimes found, and what they mean for claims of inherent privacy.}
\end{abstract}

\begin{IEEEkeywords}
federated learning, gradient inversion, variational quantum circuits, privacy
\end{IEEEkeywords}

\section{Introduction}
\fillin{Federated learning sends gradients, not data~\cite{mcmahan2017}; gradient matching inverts
them on classical networks~\cite{zhu2019}. For VQCs the literature disagrees: inherent
privacy~\cite{kumar2023}, numerical inversion once overparameterized~\cite{papadopoulos2025}, and
the trainability-privacy link~\cite{heredge2025}. Contributions.}

\section{Background and Related Work}
\subsection{Gradient Inversion in Federated Learning}
\fillin{Threat, gradient matching, why it works on classical models.}
\subsection{Variational Quantum Circuits and Data Re-uploading}
\fillin{Angle encoding, re-uploading~\cite{perezsalinas2020}, trainable blocks.}
\subsection{Competing Claims for VQCs}
\fillin{Kumar et al., Papadopoulos et al., Heredge et al., and where each should hold.}

\section{Threat Model}
\fillin{Honest-but-curious server that knows the circuit and its parameters and observes one
client's gradient.}

\section{Method}
\subsection{Circuit}
\fillin{RY angle encoding repeated $r$ times, $\ell$ trainable RY RZ blocks with a CNOT chain,
output the mean of $\braket{Z}$ over all qubits. Four qubits, then eight.}
\subsection{Attack}
\fillin{Gradient matching with L-BFGS-B, random restarts, differentiating through the simulator.}
\subsection{Outcomes}
\fillin{Recovered (distance modulo $2\pi$), ambiguous (a different input with the same gradient),
or stuck. Rates with ranges over seeds; effective parameters counted as those with nonzero
gradient.}

\section{Experiments}
\subsection{RQ1: Encoding Versus Depth}
\fillin{Fixed parameter count, more re-uploading layers versus a deeper trainable block.}
\subsection{RQ2: Client Conditions}
\fillin{Batch size, unknown label optimized jointly, trained versus random parameters.}
\subsection{RQ3: Device Noise}
\fillin{Shot noise, then fake-backend noise models. Does noise defend, and at what cost to
trainability?}

\section{Results}
\fillin{Privacy map: recovery rate over parameters and encoding repetitions, regimes marked.}

\section{Discussion and Limitations}
\fillin{What the regimes say about each claim; simulator scale; threat-model scope.}

\section{Conclusion}
\fillin{Answer to the title question.}

\bibliographystyle{IEEEtran}
\bibliography{references}

\end{document}
